\documentclass[11pt]{article}
\usepackage[T1]{fontenc}
\usepackage{lmodern,amsmath,amssymb,amsthm}
\usepackage[margin=1in]{geometry}
\usepackage[hidelinks]{hyperref}
\newtheorem{theorem}{Theorem}
\newtheorem{proposition}[theorem]{Proposition}
\newcommand{\Q}{\mathbb Q}
\newcommand{\Gal}{\operatorname{Gal}}
\newcommand{\disc}{\operatorname{disc}}
\newcommand{\rank}{\operatorname{rank}}
\title{A trace-zero point in the field of a half of a rational point}
\author{\small \href{mailto:admin@cosmosgeek.org}{admin@cosmosgeek.org}\qquad \href{mailto:b321014@iiit-bh.ac.in}{b321014@iiit-bh.ac.in}}
\date{Working paper for specialist review}
\begin{document}
\maketitle
\begin{abstract}
For an elliptic family over the curve $Y^2=X^3+8$ we exhibit two
independent rational sections and determine the rational function-field rank.
Adjoining a half of the first section produces a quartic field. An explicit
sum on the elliptic curve gives an additional nontorsion point of trace
zero in that field and a copy of the three-dimensional standard
representation in its normal closure. The construction specializes on
infinitely many rational fibres with Galois group $S_4$. A double cover of
the parameter curve gives a third rational section and a matching identity
between the relevant $2$-descent classes.
\end{abstract}

\section{The family}

Use the smooth projective parameter curve $D:Y^2=X^3+8$ and set
\[
 n=-3X/8,\qquad f(z)=z^3-nz+1/4,\qquad E_n:y^2=f(x).
\]
We restrict to finite parameters with $Y\ne0$. The cubic discriminant is
$4n^3-27/16=-27Y^2/128$, and the $j$-invariant is nonconstant.
Consider the two sections
\[
 P=(0,1/2),\qquad Q=((X-2)/4,Y/8).
\]
Let $D'$ be the smooth projective normalization obtained by adjoining
$s$ with $X+2=3s^2$. On $D'$ there is a third section $U=(1/2,3s/4)$.

\begin{theorem}\label{th:main}
The rank of $E_n$ over $\Q(D)$ is two, and $P,Q$ are independent there.
Over $\Q(D')$, the points $P,Q,U$ are independent. There are infinitely
many rational parameters on $D'$ giving pairwise distinct $j$-invariants
with the following properties. All three specialized rational points are
independent. For any half $H$ of $P$, the field $K=\Q(H)$ is quartic
with $S_4$ normal closure $L$, and there is an explicit nontorsion
$S_H\in E_n(K)$ of trace zero. In particular,
\[
 \rank E_n(K)\ge\rank E_n(\Q)+1,\qquad
 \rank E_n(L)\ge\rank E_n(\Q)+3.
\]
\end{theorem}

Quartic rank growth for a fixed elliptic curve has been studied in
families of $S_4$ extensions \cite{LOT}. Here the curve varies and the
extension is prescribed by a marked half-point. The specialization
argument gives infinitely many curves and makes no assertion about
the number of distinct quartic fields. Section~\ref{sec:descent}
gives an explicit relation on $2$-descent classes.

\section{Rational sections and generic rank}

At $T=(X,Y)=(1,3)$, the change of coordinates $(x,y)\mapsto(4x,8y)$
gives the integral equation $y^2=x^3+6x+16$. At $s=1$ the three sections
specialize to $(0,4),(-1,3),(2,6)$. Their doubles have respective
$x$-coordinates $9/16,17/4,-7/4$. Nagell--Lutz shows that each
specialization has infinite order. In particular the three sections are
generically nontorsion. Their equations follow by substitution, using
$f(1/2)=3(X+2)/16$. The involution $Y\mapsto-Y$ separates $Q$ from $P$.
The squareclasses $X^3+8$ and $(X+2)/3$ are independent in
$\overline{\Q}(X)^\times/\overline{\Q}(X)^{\times2}$, since their
product has simple zeros at the roots of $X^2-2X+4$. Thus changing the
signs of $s$ and $Y$ independently separates $P,Q,U$ into three
different characters. This proves independence over the respective
function fields.

For the upper bound, pass first to $\overline{\Q}(X)$. At the three
simple finite zeros of $X^3+8$ the fibre is $I_1$. At infinity, using
the parameter $t=1/X$ and replacing $(x,y)$ by $(t^2x,t^3y)$ gives
valuations $3,6,9$ for $a_4,a_6,\Delta$, respectively. The fibre is
$III^*$. The resulting rational elliptic surface has Picard number $10$,
so Shioda--Tate gives geometric Mordell--Weil rank $10-2-7=1$
\cite{SS}. The quadratic twist by $X^3+8$ has the three fibres $I_1^*$
and one $III$, by the quadratic twisting rules. Its fibre Euler numbers
sum to $24$, so it is an elliptic K3 surface and has Picard number at
most $20$. Its geometric Mordell--Weil rank is consequently at most
$20-2-3\cdot5-1=2$.

To determine the rational part of this twist, choose
$\zeta^2+\zeta+1=0$. The three points
\[
 Q_j=\left(\frac{\zeta^jX-2\zeta^{-j}}4,\frac{Y}{8}\right),
 \qquad j=0,1,2,
\]
all lie on one horizontal line, so they sum to zero. They belong to the
anti-invariant part for the cover $D\to\mathbb P^1_X$, which identifies
over $\Q$-coefficients with the twist just discussed. The order-three
operator on sections taking $(x(X,Y),y(X,Y))$ to
$(\zeta\,x(\zeta X,Y),y(\zeta X,Y))$ cycles these points. Their
$\Q$-linear span in the group tensored with $\Q$ has dimension two:
it is nonzero, and on a one-dimensional span the operator would be
trivial, contradicting $Q_0+Q_1+Q_2=0$ and nontorsion. They therefore
span the whole geometric anti-invariant part. Complex conjugation fixes
$Q_0$ and exchanges $Q_1,Q_2$, and its fixed line is spanned by $Q_0$.
The invariant part for $D\to\mathbb P^1_X$ already has rank one, generated
up to finite index by $P$. This proves the claimed generic rank over
$\Q(D)$.

\section{The marked quartic and its extra point}

The four tangent lines at the halves of $P$ pass through $-P$. Substitution
of $y=tx-1/2$ gives $x(x^2-t^2x+t-n)=0$. The tangency equation and
point are thus
\[
 t^4-4t+4n=0,\qquad H=(t^2/2,(t^3-1)/2).
\]
The tangent slope is recovered from $H$, so $K=\Q(t)$ on a smooth
rational fibre. Let $\theta_i$ be the three roots of $f$ and choose
$r_i^2=-\theta_i$ with $r_1r_2r_3=1/2$. The four signed sums
$t=\sum_i r_i$ with this product are the four roots of the tangent
equation. Indeed $\sum r_i^2=0$ and
$(\sum_{i<j} r_i r_j)^2=-n+t$. The sums are distinct because $f$ is
separable with nonzero roots. Each $r_i$ is half the sum of $t$ and
the root obtained by negating the other two signs. This identifies the
quartic splitting field with that of the signed triples. These are the
usual square roots appearing in explicit halving \cite{BZ}.

For the choice belonging to $H$, set
\begin{equation}\label{eq:ri}
 R_i=\left(-\frac{\theta_i}{2},\frac{9Y r_i}{32f'(\theta_i)}\right),
 \qquad S_H=R_1+R_2+R_3.
\end{equation}
The second operation is elliptic-curve addition. For the equation in
\eqref{eq:ri}, one may reduce both sides of
\[
 1024 f'(\theta)^2 f(-\theta/2)=-81Y^2\theta
\]
modulo $\theta^3-n\theta+1/4$ and use $Y^2=8-512n^3/27$.
The stabilizer of the signed sum $t$ permutes the three terms of
\eqref{eq:ri}, so $S_H$ is defined over the marked quartic field.
Summing $S_H$ over the four sign choices gives zero: each $R_i$ occurs
twice with each sign.

Here is an independent certificate of generic nontorsion of each $R_i$.
On adjoining $r=r_i$ one has $n=r^4-1/(4r^2)$. At $r=1/2$ this
gives $n=-15/16$, $X=5/2$ and $Y=\pm3\sqrt{42}/4$. The point $R_i$
specializes to $(1/8,\pm3\sqrt{42}/32)$ on a smooth fibre. These
values occur on each component in use: the nonconstant $r$ on the
smooth projective parameter cover takes the value $1/2$, and either
choice for $Y$ gives the same nontorsion test. On the quadratic twist
by $d=42$, with the further coordinate scaling $(x,y)\mapsto(4x,8y)$,
this becomes the rational point $(21,\pm1323)$ on
\[
 y^2=x^3+26460x+1185408 .
\]
Its double has $x=273/4$, so Nagell--Lutz proves nontorsion.
Silverman's specialization theorem \cite{Silverman} applies on
the finite parameter cover where the $R_i$ are defined. The
canonical height of a nontorsion section grows with the base height,
as the $j$-invariant is nonconstant. Points above rational base
parameters have bounded degree, so Northcott's theorem leaves only
finitely many exceptional specializations.

Assume now that the tangent quartic has Galois group $S_4$. Its normal
Klein four subgroup acts by the even sign changes. If $\sigma$ negates
the two signs other than $r_i$, then $S_H+\sigma S_H=2R_i$. Consequently
$S_H$ has infinite order whenever $R_i$ does. A nontorsion point of
trace zero is independent of $E_n(\Q)$ in $E_n(K)\otimes\Q$. Moreover
the four conjugates define a nonzero $S_4$-equivariant map from the
permutation module to $E_n(L)\otimes\Q$ that kills the constant vector.
The augmentation summand of this module is irreducible of dimension
three. Its image therefore gives the claimed three-dimensional rank
gain in the normal closure.

\section{Infinitely many simultaneous specializations}

The function $X+2$ has divisor $2[(-2,0)]-2[O]$ on $D$. It is not a
square over $\overline{\Q}(D)$ since $[(-2,0)]-[O]$ is nonprincipal.
Thus $D'\to D$ is a connected unramified double cover. Riemann--Hurwitz
gives genus one, and $(1,3,1)$ supplies a rational point on $D'$.
With a choice of origins the map differs from a rational isogeny by
translation. The point $T=(1,3)$ has infinite order on $D$, as its
double has $x=-7/4$. The dual isogeny then shows $D'(\Q)$ has positive
rank.

At $(X,Y)=(1,3)$ the tangent quartic is $g(t)=t^4-4t-3/2$.
Modulo $5$ this has the single root $3$, and the complementary cubic
is irreducible (it has no root and the root $3$ is simple).
Modulo $19$ the quartic has no factor of degree one or two.
For a short check, $t^{361}-t$ has residue
$4t^3+13t^2+8t+7$ modulo $g$ at $19$; the Euclidean algorithm with $g$
then gives successive nonzero remainders
$5t^2+15t+3,14t+14,12$. The factorization types give a three-cycle
and a four-cycle in the Galois group. A subgroup of $S_4$ containing
both has order divisible by $12$ and contains an odd permutation,
hence is $S_4$.

On $D'(\Q_5)$ and $D'(\Q_{19})$ one may retain these congruences on open
neighborhoods of $(1,3,1)$. The intersection with $D'(\Q)$ contains
a coset of a finite-index subgroup: a sufficiently small open subgroup
of each compact local elliptic group suffices after translation.
This gives infinitely many rational parameters. Specialization preserves
the independence of $P,Q,U$ outside finitely many of them
\cite{Silverman}, and the preceding argument handles $R_i$. Both the map
to the $X$-line and the $j$-map are nonconstant and have finite fibres.
This finishes the proof of Theorem~\ref{th:main}.

\section{The matching descent classes}\label{sec:descent}

The third section on $D'$ can also be seen in the cubic algebra of the
original family. Work over $k=\Q(D)$ with $B=k[\theta]=k[z]/f(z)$ and the
canonical identifications of the two-torsion of quadratic twists.
The Kummer class of $(x,y)$ on $E_n$ is represented by $x-\theta$
in the norm-square kernel of $B^\times/B^{\times2}$. On the $d$-twist,
a point mapped to $(x,y)$ over $k(\sqrt d)$ is represented by
$d(x-\theta)$. These descriptions follow by splitting the exact sequence
of the norm on the restriction of scalars of $\mu_2$ (the degree is
odd) and using the Weil-pairing functions $x-\theta_i$.

The single term $R$ of \eqref{eq:ri} over $B(\sqrt{-\theta})$ is
anti-invariant under the sign change and hence gives a $B$-point on the
$(-\theta)$-twist. Corestrict its Kummer class from $B$ to $k$.
Under the preceding description, corestriction is the norm from
$B\otimes_k B$ to $B$ in the first coordinate. A representative
of the result is
\begin{align*}
\prod_i(-\theta_i)(-\theta_i/2-\theta)
 &= -\frac{1}{32}\prod_i(\theta_i+2\theta)\\
 &= -\frac{3}{32}\theta f'(\theta)
 \ \equiv\ -6\theta f'(\theta)\pmod {B^{\times2}}.
\end{align*}
Here $-f(-2\theta)=3\theta f'(\theta)$.

Put $d=3(X+2)$ and take the $d$-twist point corresponding to
$(1/2,\sqrt d/4)$ on $E_n$. Using $n=\theta^2+1/(4\theta)$ yields
the exact identity
\[
 (x_Q-\theta)d(1/2-\theta)=
 (-6\theta f'(\theta))
 \left(\frac{(\theta+1)(2\theta-1)}{3\theta}\right)^2 .
\]
Thus the corestriction just computed is the sum of the class of $Q$
and this twist class. On $D'$ the latter becomes the class of $U$.
This identity is at the level of $2$-descent. The point $S_H$ itself
retains trace zero over the quartic field.

\section{Central vanishing}

The established analytic-rank-zero and analytic-rank-one theorems
imply that rational rank at least two forces analytic rank at least
two \cite{Clay}. Therefore the specializations in
Theorem~\ref{th:main} have $L(E_n,1)=L'(E_n,1)=0$.
With negative root number the order is at least three. No recurrence
of either root number on the specified coset is used.

The odd-prime part of the sign has a simple independent description.
Write $X=A/B^2$ and $Y=C/B^3$ in lowest terms with $B>0$.
An integral model has coefficients $a_4=6AB^2$ and $a_6=16B^6$.
For $p\ge5$ dividing $C$, its $c_4$ is a unit and its discriminant
vanishes mod $p$, since it is $-27\cdot 2^9 B^6 C^2$. The reduction
is multiplicative, split precisely when $-c_6$ is a square mod $p$,
or equivalently $(6/p)=1$. For $p\ge5$ dividing $B$, the integers
$A,C$ are units at $p$. If $v_p(B)$ is even, the reduction is good
after scaling. If it is odd, the curve is a ramified quadratic twist
of a curve with good reduction: after scaling away the even part of
$v_p(B)$ the valuations of $a_4,a_6,\Delta$ are $2,\ge3,6$.
Its local sign is $(-1/p)$ by the local root-number formula
\cite{Root}. All remaining primes $p\ge5$ are good. Hence, if $B'$
is the positive part of $B$ prime to $6$, the global sign is
\[
 -w_2(E_n)w_3(E_n)
 \left(\frac{-1}{B'}\right)
 (-1)^{\#\{p\ge5:p\mid C,\ (6/p)=1\}},
\]
with the Jacobi symbol in the middle. Determining a central second
derivative or proving exact analytic rank requires further input.

\begin{thebibliography}{9}
\bibitem{BZ} B.~M. Bekker and Y.~G. Zarhin,
\emph{The divisibility by 2 of rational points on elliptic curves}.
\href{https://arxiv.org/abs/1702.02255}{arXiv:1702.02255}.
\bibitem{SS} M. Sch\"utt and T. Shioda, \emph{Elliptic surfaces},
especially the discussions of quadratic twists and the Shioda--Tate formula.
\href{https://arxiv.org/abs/0907.0298}{arXiv:0907.0298}.
\bibitem{Silverman} J.~H. Silverman,
\emph{Heights and the specialization map for families of abelian varieties},
J. reine angew. Math. \textbf{342} (1983), 197--211.
\href{https://doi.org/10.1515/crll.1983.342.197}{doi:10.1515/crll.1983.342.197}.
\bibitem{LOT} R.~J. Lemke Oliver and F. Thorne,
\emph{Rank growth of elliptic curves in nonabelian extensions}.
\href{https://arxiv.org/abs/1810.04018}{arXiv:1810.04018}.
\bibitem{Root} D. Rohrlich, \emph{Elliptic curves and the Weil--Deligne group},
in \emph{Elliptic Curves and Related Topics}, CRM Proc. Lecture Notes
\textbf{4} (1994), 125--157.
\bibitem{Clay} A. Wiles, \emph{The Birch and Swinnerton-Dyer conjecture}.
\href{https://www.claymath.org/wp-content/uploads/2022/05/birchswin.pdf}{Clay problem description}.
\end{thebibliography}
\end{document}
